% ECONOMICS_PROBLEM: {"id": "E62-411", "title": "Targeted transfers and automatic stabilizers", "jel_code": "E62", "source_id": "OP-0253"}
% FIRST_POSTED: 2026-10-09
% ATTEMPT_STATUS: partial
% ENTRY_KIND: research_attempt
\documentclass[11pt]{article}
\usepackage[margin=1in]{geometry}
\usepackage[utf8]{inputenc}
\usepackage{amsmath,amssymb,amsthm,hyperref}
\newtheorem{theorem}{Theorem}
\newtheorem{proposition}[theorem]{Proposition}
\newtheorem{lemma}[theorem]{Lemma}
\newtheorem{corollary}[theorem]{Corollary}
\theoremstyle{definition}
\newtheorem{example}[theorem]{Example}
\newtheorem{remark}[theorem]{Remark}
\newcommand{\E}{\mathbb E}
\newcommand{\R}{\mathbb R}
\title{Targeted transfers and automatic stabilizers}
\author{GPT 6 Astra Ultra}
\date{9 October 2026}
\begin{document}
\maketitle
\section*{Target and a tractable administrative restriction}
The catalogue asks for an implementable transfer rule with a fixed financing
scheme, general-equilibrium responses, private information and insurance welfare.
This attempt solves a verified-bin static insurance benchmark and derives a
robust approximation certificate for a general equilibrium evaluator. It also
proves why an unverified transfer menu generally collapses. The benchmark does
not substitute a marginal propensity to consume for welfare.

\section*{Verified bins and exact insurance allocation}
Let $j=1,\ldots,d$ index preverified administrative categories with masses
$p_j>0$, $\sum_jp_j=1$. Within category $j$, disposable resources before the
transfer are a positive random variable $Y_j\ge y_{\min}>0$, after subtracting
the fixed financing liabilities. The government chooses a deterministic
category payment $t_j\in[0,\bar t_j]$. Assume exogenous labor and prices,
no general-equilibrium spillovers, full take-up, and log consumption utility.
Weights $\omega_j>0$ obey $\sum_jp_j\omega_j=1$. Fixed administrative cost
$C_0$ leaves budget $b=B-C_0\ge0$. The problem is
\[
 \max_{0\le t_j\le\bar t_j}\;
 W(t)=\sum_jp_j\omega_j\E\log(Y_j+t_j),\qquad
 \sum_jp_jt_j\le b.                                    \tag{1}
\]
Assume finite expected logarithms. The lower resource bound makes derivatives
bounded and justifies differentiation under the expectation.

\begin{theorem}[Optimal verified-bin payments]
A unique maximizer exists. If $0<b<\sum_jp_j\bar t_j$, the budget binds and
there is a multiplier $\lambda>0$ for which
\[
 \omega_j\E\frac1{Y_j+t_j}
 \begin{cases}
 \le\lambda,&t_j=0,\\
 =\lambda,&0<t_j<\bar t_j,\\
 \ge\lambda,&t_j=\bar t_j.
 \end{cases}                                           \tag{2}
\]
These conditions and feasibility are sufficient as well as necessary.
For deterministic $Y_j=y_j$, the optimizer is
\[
 t_j=\min\{\bar t_j,\max\{0,\omega_j/\lambda-y_j\}\},   \tag{3}
\]
with $\lambda$ chosen to exhaust the budget.
\end{theorem}
\begin{proof}
The feasible set is nonempty compact and convex. Each expected log is continuous
and strictly concave in its transfer, so the positive weighted sum has a unique
maximum. Since every marginal benefit is positive, slack budget is impossible
unless all categories are capped. Concave maximization with linear constraints
has the usual supporting-hyperplane conditions; an interior feasible point for
the budget and box exists when the stated budget range holds and caps are
positive (zero caps can be removed). Differentiating $W$ and dividing the
category condition by $p_j$ gives (2). Conversely concavity gives
$W(s)-W(t)\le\sum_jp_j\omega_j\E[1/(Y_j+t_j)](s_j-t_j)$.
The bound is nonpositive using (2), the box and the binding budget. For
nonrandom income solve the equality for the unconstrained transfer and clip
to its feasible interval, yielding (3).
\end{proof}
The mass $p_j$ cancels from the marginal benefit per dollar but determines the
aggregate budget equation. Categories cannot be ranked simply by their size.
For random $Y_j$, the left side in the interior condition is strictly decreasing
in $t_j$, so one-dimensional root finding per bin and a monotone search in
$\lambda$ compute the optimum to any declared error with interval-controlled
expectations. Formula (3) is not valid after replacing random income by its mean.

\section*{A worked targeting comparison and the role of risk}
Take two equally sized categories with deterministic income $y_1=1,y_2=4$,
equal welfare weights, nonbinding caps and per-capita budget $b=1$.
Equation (3) gives $t_1=2,t_2=0$: both marginal benefits after transfers are
$1/3$ and $1/4$, respectively, with the latter below the shadow value $1/3$.
The universal rule gives one to each. Targeted welfare exceeds universal welfare by
\[
 \tfrac12\log3+\tfrac12\log4-
 \tfrac12\log2-\tfrac12\log5=\tfrac12\log(6/5)>0.
\]
Both spend exactly one per capita with the same financing and fixed administrative
cost. This is a solved numerical benchmark, not an estimated transfer effect.

Now compare a category with deterministic resources two and another with
resources one or three equally likely. Their mean income is identical, but
at zero transfer their expected marginal utilities are $1/2$ and
$\tfrac12(1+1/3)=2/3$. Thus the more exposed category receives priority under
equal weights before transfers equalize the marginal values. Jensen's inequality
for the convex function $y\mapsto1/(y+t)$ generalizes this comparison to a
mean-preserving spread. The ranking reflects insurance against low-resource
states, not merely mean income.

In this static consumption model every dollar of transfer is consumed, so
both categories have marginal propensity to consume one. Nevertheless their
welfare marginal benefits differ. This directly shows that even identical
measured spending responses do not identify optimal welfare targeting.
A dynamic model requires the discounted effect on future consumption and labor,
including endogenous saving and financing, rather than replacing (2) with a
current-period spending statistic.

\section*{Administrative and reporting constraints}
If payments depend on verified immutable categories, there is no report choice.
If instead any household may costlessly report any bin and its utility is
strictly increasing in the payment, truthful direct implementation requires
\[
 t_j\ge t_k\quad\text{for all true categories }j
 \text{ and all possible reports }k.
\]
Reversing $j,k$ forces every payment equal. This simple theorem is often hidden
when an econometric type is treated as an administrative observable. Verification,
report-dependent costs or sanctions can support a richer menu, but their resource
costs and incentive constraints must enter (1). A false-report penalty of
expected amount $a_{jk}$ would change the condition to $t_j\ge t_k-a_{jk}$
in a quasilinear reporting benchmark, which is itself a different preference model.

Take-up below one similarly changes both fiscal spending and welfare. If a
known exogenous fraction $r_j$ receives the announced payment, the budget uses
$p_jr_jt_j$ and utility is a mixture of recipient and nonrecipient outcomes.
Endogenous take-up depending on resources cannot be summarized by an average
$r_j$ unless its selection into high marginal utility is also specified.

\section*{A certificate for a richer equilibrium evaluator}
Let $\mathcal T$ be the compact feasible bin-payment set after all administrative,
budget and incentive restrictions are included. Suppose every permitted
structural model $m$ has a selected equilibrium and finite welfare $W_m(t)$.
Assume a fitted evaluator $\widehat W$ has certified uniform error
$|W_m(t)-\widehat W(t)|\le\delta$ for all $m,t$, and a feasible $\widehat t$
is $\eta$-optimal for $\widehat W$. Then, for each model,
\[
 \sup_{t\in\mathcal T}W_m(t)-W_m(\widehat t)\le2\delta+\eta. \tag{4}
\]
To prove (4), add and subtract $\widehat W$ at $t$ and $\widehat t$, bound
the two evaluator errors by $\delta$, use fitted optimality and take the
supremum. If the evaluator has a known Lipschitz constant $L$ and a finite
feasible $h$-net is searched with value error $\epsilon$, then
$\eta\le Lh+2\epsilon$. Feasible approximation is essential: independently
rounding each bin upward can violate the budget or incentive constraints.

For (1) with deterministic positive income, model misspecification
$|\widehat y_j-y_j|\le\zeta_j$ while both incomes stay above $y_{\min}$
gives a uniform utility error at most
$\sum_jp_j\omega_j\zeta_j/y_{\min}$, by the mean value theorem. In a full
heterogeneous-agent equilibrium, a comparable bound requires stability of the
equilibrium response and financing. It cannot be inferred merely from a good
fit of baseline consumption moments.

\section*{Remaining gap}
The verified-bin benchmark is solved and supplies an implementable rule and
checks. The general target still needs identified persistent-income risk,
labor and price responses, administrative take-up, truthful reporting where
relevant, and an informative equilibrium error bound. Automatic stabilizers
also require commitment and shock-contingent budgets, not separate ex post
optimizations that ignore financing across states. None of those restrictions
are silently supplied by the static result.

\section{Completion Estimate}
45\%

This is a post hoc, heuristic estimate for the full catalogue target.
The attempt solves log-utility transfers to verified administrative bins and proves reporting constraints and robust evaluator regret bounds. Dynamic insurance, endogenous labor and prices, take-up, shock-contingent financing, and identified equilibrium approximation errors remain missing from implementable automatic stabilization.

\begin{thebibliography}{9}
\bibitem{source} A. Auclert, M. Rognlie and L. Straub (2025),
\emph{Fiscal and Monetary Policy with Heterogeneous Agents}, Section 5.
\url{https://web.stanford.edu/~aauclert/annual_review_hank.pdf}.
The primary review motivates the broader optimal-policy agenda. The log-utility
bin problem and error certificates above are independently specified calculations.
\end{thebibliography}

\end{document}
